% TODO checklist:
% - [x] Inspected src/background.py (327 lines - complete background framework)
% - [x] Inspected src/services/health.py (HealthService)
% - [x] Inspected src/services/archive.py (ArchiveService)
% - [x] Inspected src/services/service_metrics.py (ServiceMetrics)
% - [x] Analyzed Q&A.md Q40 for scheduled task categories

\section{Background Framework}
\label{sec:background-framework}

Beyond the job queue for user-initiated tasks, the Cineca Agentic Platform requires periodic maintenance operations: health monitoring, data backups, cache cleanup, and index maintenance. The background framework provides a unified scheduler for these operations, built on APScheduler with comprehensive metrics integration.

This section describes the background framework architecture, scheduled task categories, and integration with the FastAPI application lifecycle.

\subsection{Framework Overview}

The background framework in \texttt{src/background.py} coordinates periodic tasks using the APScheduler library with the AsyncIO scheduler backend. Key design goals include:

\begin{itemize}
    \item \textbf{Configurable Scheduling}: Tasks can be scheduled via interval triggers (e.g., every 30 seconds) or cron expressions (e.g., daily at 02:30 UTC).
    
    \item \textbf{Graceful Lifecycle}: Tasks are started with the FastAPI application and stopped cleanly on shutdown.
    
    \item \textbf{Metrics Integration}: All task executions are instrumented with duration and status metrics.
    
    \item \textbf{Failure Isolation}: Task failures are logged but do not affect other scheduled tasks.
    
    \item \textbf{Service Composition}: Tasks delegate to specialized services (HealthService, ArchiveService) for actual operations.
\end{itemize}

\subsection{Configuration Model}

The \texttt{BackgroundConfig} dataclass captures all scheduling parameters:

\begin{lstlisting}[language=Python, caption={Background framework configuration.}]
@dataclass
class BackgroundConfig:
    enabled: bool = getattr(settings, "BACKGROUND_ENABLED", True)

    # Health probe interval (seconds)
    health_enabled: bool = getattr(settings, "BACKGROUND_HEALTH_ENABLED", True)
    health_interval_seconds: int = getattr(
        settings, "BACKGROUND_HEALTH_INTERVAL_SECONDS", 30
    )

    # Backups via cron (crontab format, UTC)
    backup_enabled: bool = getattr(settings, "BACKGROUND_BACKUPS_ENABLED", False)
    backup_cron: str = getattr(
        settings, "BACKGROUND_BACKUPS_CRON", "30 2 * * *"
    )  # 02:30 UTC daily

    # Cleanup via cron (crontab format, UTC)
    cleanup_enabled: bool = getattr(settings, "BACKGROUND_CLEANUP_ENABLED", False)
    cleanup_cron: str = getattr(
        settings, "BACKGROUND_CLEANUP_CRON", "15 3 * * 0"
    )  # 03:15 UTC Sundays

    # Redis job store index cleanup
    redis_cleanup_enabled: bool = getattr(
        settings, "BACKGROUND_REDIS_CLEANUP_ENABLED", True
    )
    redis_cleanup_interval_seconds: int = getattr(
        settings, "BACKGROUND_REDIS_CLEANUP_INTERVAL", 3600
    )  # 1 hour
    redis_cleanup_batch_size: int = getattr(
        settings, "BACKGROUND_REDIS_CLEANUP_BATCH_SIZE", 500
    )
\end{lstlisting}

\subsection{BackgroundManager Class}

The \texttt{BackgroundManager} class orchestrates task scheduling and execution:

\begin{lstlisting}[language=Python, caption={BackgroundManager initialization.}]
class BackgroundManager:
    """Coordinates background jobs with metrics and logging."""

    def __init__(
        self,
        *,
        health: HealthService | None = None,
        archive: ArchiveService | None = None,
        metrics: ServiceMetrics | None = None,
        config: BackgroundConfig | None = None,
    ) -> None:
        self.config = config or BackgroundConfig()
        self.scheduler: AsyncIOScheduler | None = None
        self.health = health or HealthService()
        self.archive = archive or ArchiveService()
        self.metrics = metrics or ServiceMetrics()
        self._started = False
\end{lstlisting}

\subsubsection{Lifecycle Management}

The manager provides async start/stop methods for integration with application lifecycle:

\begin{lstlisting}[language=Python, caption={Background manager lifecycle methods.}]
async def start(self) -> None:
    """Start the underlying scheduler and register jobs."""
    if not self.config.enabled:
        log.info("background.disabled")
        return

    if self._started:
        log.debug("background.already_started")
        return

    self.scheduler = AsyncIOScheduler(timezone="UTC")
    self._register_jobs()
    self.scheduler.start()
    self._started = True
    log.info(
        "background.started",
        health_enabled=self.config.health_enabled,
        backups_enabled=self.config.backup_enabled,
        cleanup_enabled=self.config.cleanup_enabled,
    )

async def stop(self) -> None:
    """Stop scheduler and wait for running jobs to finish."""
    if not self._started or not self.scheduler:
        return
    try:
        self.scheduler.shutdown(wait=True)
        log.info("background.stopped")
    finally:
        self._started = False
        self.scheduler = None
\end{lstlisting}

\subsubsection{Job Registration}

Tasks are registered with APScheduler using interval or cron triggers:

\begin{lstlisting}[language=Python, caption={Task registration with APScheduler.}]
def _register_jobs(self) -> None:
    assert self.scheduler is not None

    if self.config.health_enabled:
        self.scheduler.add_job(
            self._wrap_job(self._job_health, "health"),
            IntervalTrigger(seconds=max(1, int(self.config.health_interval_seconds))),
            id="background.health",
            replace_existing=True,
            coalesce=True,  # Combine missed executions
            max_instances=1,  # Prevent concurrent runs
        )

    if self.config.backup_enabled:
        with contextlib.suppress(ValueError):
            self.scheduler.add_job(
                self._wrap_job(self._job_backup, "backup"),
                CronTrigger.from_crontab(self.config.backup_cron),
                id="background.backup",
                replace_existing=True,
                coalesce=True,
                max_instances=1,
            )

    if self.config.cleanup_enabled:
        with contextlib.suppress(ValueError):
            self.scheduler.add_job(
                self._wrap_job(self._job_cleanup, "cleanup"),
                CronTrigger.from_crontab(self.config.cleanup_cron),
                id="background.cleanup",
                replace_existing=True,
                coalesce=True,
                max_instances=1,
            )

    # Redis cleanup (only if using Redis backend)
    if (self.config.redis_cleanup_enabled and 
        settings.JOB_STORE_BACKEND == "redis"):
        self.scheduler.add_job(
            self._wrap_job(self._job_redis_cleanup, "redis_cleanup"),
            IntervalTrigger(
                seconds=max(60, int(self.config.redis_cleanup_interval_seconds))
            ),
            id="background.redis_cleanup",
            replace_existing=True,
            coalesce=True,
            max_instances=1,
        )
\end{lstlisting}

\subsection{Job Wrapper with Metrics}

All scheduled tasks are wrapped in a metrics-recording decorator:

\begin{lstlisting}[language=Python, caption={Job wrapper for metrics and error handling.}]
def _wrap_job(
    self, func: Callable[[], Awaitable[Any]], job_name: str
) -> Callable[[], Awaitable[None]]:
    async def runner() -> None:
        start = time.perf_counter()
        status = "ok"
        try:
            await func()
        except Exception as e:
            status = "error"
            log.warning("background.job.error", job=job_name, err=str(e))
        finally:
            dur = time.perf_counter() - start
            # Record metrics if available
            with contextlib.suppress(Exception):
                if hasattr(self.metrics, "record_bg_job"):
                    self.metrics.record_bg_job(
                        job_name, status=status, duration_seconds=dur
                    )
            log.debug(
                "background.job.done",
                job=job_name, status=status, duration=f"{dur:.3f}s"
            )
    return runner
\end{lstlisting}

\subsection{Scheduled Task Categories}
\label{subsec:scheduled-task-categories}

The background framework includes four categories of scheduled tasks:

\subsubsection{Health Monitoring Tasks}

The health job performs periodic checks of all system components:

\begin{lstlisting}[language=Python, caption={Health monitoring background task.}]
async def _job_health(self) -> None:
    """Periodic health check sweep."""
    res = await self.health.check()
    if not res.ok:
        log.warning("background.health.unhealthy", error=res.error)
    else:
        log.info("background.health.ok")
\end{lstlisting}

Components checked include:
\begin{itemize}
    \item PostgreSQL connectivity and query capability
    \item Redis ping latency and memory usage
    \item Memgraph connection status (if configured)
    \item LLM provider health probes
\end{itemize}

\subsubsection{Backup Tasks}

The backup job delegates to ArchiveService for creating snapshots:

\begin{lstlisting}[language=Python, caption={Backup background task.}]
async def _job_backup(self) -> None:
    """Create a snapshot/backup via ArchiveService."""
    fn = None
    # Support multiple method names for flexibility
    for candidate in ("backup", "run_backup", "create_backup"):
        if hasattr(self.archive, candidate):
            fn = getattr(self.archive, candidate)
            break
    
    if fn is None:
        log.warning("background.backup.skipped", reason="no_method")
        return

    maybe_awaitable = fn()
    if asyncio.iscoroutine(maybe_awaitable):
        await maybe_awaitable
    log.info("background.backup.completed")
\end{lstlisting}

Backup operations include:
\begin{itemize}
    \item Memgraph graph snapshots (JSON + GZIP compression)
    \item Redis RDB/AOF backup triggers
    \item Manifest generation with SHA-256 checksums
\end{itemize}

\subsubsection{Cleanup Tasks}

The cleanup job handles age-based data pruning:

\begin{lstlisting}[language=Python, caption={Cleanup background task.}]
async def _job_cleanup(self) -> None:
    """Run retention/cleanup via ArchiveService."""
    fn = None
    for candidate in ("cleanup", "run_cleanup", "prune"):
        if hasattr(self.archive, candidate):
            fn = getattr(self.archive, candidate)
            break
    
    if fn is None:
        log.info("background.cleanup.skipped", reason="no_method")
        return

    maybe_awaitable = fn()
    if asyncio.iscoroutine(maybe_awaitable):
        await maybe_awaitable
    log.info("background.cleanup.completed")
\end{lstlisting}

Cleanup operations include:
\begin{itemize}
    \item Expired session removal
    \item Old backup file pruning (retention policy)
    \item Stale cache entry eviction
    \item Completed job record cleanup (beyond TTL)
\end{itemize}

\subsubsection{Redis Index Cleanup}

The Redis cleanup job removes orphaned entries from ZSET indexes:

\begin{lstlisting}[language=Python, caption={Redis index cleanup task.}]
async def _job_redis_cleanup(self) -> None:
    """Clean orphaned ZSET members from Redis job store indexes."""
    if settings.JOB_STORE_BACKEND != "redis":
        return

    from db.redis_cache.async_client import get_async_redis
    from db.redis_cache.job_store import RedisJobStore

    store = RedisJobStore()
    redis = await get_async_redis()
    batch_size = self.config.redis_cleanup_batch_size
    total_removed = 0

    try:
        # 1. Clean global index
        removed = await store.cleanup_orphaned_index_members(
            "jobs:all", batch_size=batch_size
        )
        total_removed += removed

        # 2. Clean status indexes
        for status in ["queued", "running", "finished", "failed", "cancelled"]:
            index_key = f"jobs:status:{status}"
            removed = await store.cleanup_orphaned_index_members(
                index_key, batch_size=batch_size
            )
            total_removed += removed

        # 3. Clean owner indexes (scan with cursor)
        cursor = 0
        owner_indexes_cleaned = 0

        while True:
            cursor, keys = await redis.scan(
                cursor=cursor, match="jobs:owner:*", count=100
            )
            for key_bytes in keys:
                key = (key_bytes.decode("utf-8") 
                       if isinstance(key_bytes, bytes) else key_bytes)
                removed = await store.cleanup_orphaned_index_members(
                    key, batch_size=batch_size
                )
                total_removed += removed
                owner_indexes_cleaned += 1
            if cursor == 0:
                break

        log.info(
            "background.redis_cleanup.completed",
            total_orphans_removed=total_removed,
            owner_indexes_cleaned=owner_indexes_cleaned,
        )

    except Exception as e:
        log.error("background.redis_cleanup.error", err=str(e))
        raise
\end{lstlisting}

Orphaned index members occur when job HASH keys expire via TTL but their ZSET entries remain. The cleanup uses a Lua script to atomically check existence and remove orphans:

\begin{lstlisting}[language=Python, caption={Orphan cleanup Lua script (conceptual).}]
CLEANUP_ORPHANS_SCRIPT = """
-- KEYS[1] = index key (e.g., "jobs:all")
-- ARGV[1] = batch_size

local members = redis.call('ZRANGE', KEYS[1], 0, ARGV[1] - 1)
local removed = 0

for _, job_id in ipairs(members) do
    local exists = redis.call('EXISTS', 'job:' .. job_id)
    if exists == 0 then
        redis.call('ZREM', KEYS[1], job_id)
        removed = removed + 1
    end
end

return removed
"""
\end{lstlisting}

\subsection{FastAPI Lifespan Integration}

The background framework provides a context manager for FastAPI lifespan integration:

\begin{lstlisting}[language=Python, caption={FastAPI lifespan integration.}]
def build_default_manager() -> BackgroundManager:
    """Create a BackgroundManager with default services and config."""
    return BackgroundManager(
        health=HealthService(),
        archive=ArchiveService(),
        metrics=ServiceMetrics(),
        config=BackgroundConfig(),
    )


@asynccontextmanager
async def lifespan(app: FastAPI):
    """Attach a default BackgroundManager to app.state.bg."""
    manager = build_default_manager()
    app.state.bg = manager
    await manager.start()
    try:
        yield
    finally:
        await manager.stop()
\end{lstlisting}

Usage in the FastAPI application:

\begin{lstlisting}[language=Python, caption={Lifespan usage in app.py.}]
from src.background import lifespan

app = FastAPI(
    title="Cineca Agentic Platform",
    lifespan=lifespan,
    # ... other configuration
)
\end{lstlisting}

\subsection{Metrics Emission}

Background tasks emit Prometheus metrics for operational visibility:

\begin{lstlisting}[language=Python, caption={Background job metrics (conceptual).}]
background_jobs_total = Counter(
    "background_jobs_total", "Total background job executions",
    ["job", "status"]  # job: health, backup, cleanup; status: ok, error
)

background_job_duration_seconds = Histogram(
    "background_job_duration_seconds", "Background job duration",
    ["job", "status"],
    buckets=(0.01, 0.05, 0.1, 0.5, 1.0, 5.0, 10.0, 30.0, 60.0, 300.0)
)
\end{lstlisting}

\subsection{Configuration Reference}

Table~\ref{tab:background-config} summarizes the background framework configuration parameters.

\begin{table}[ht]
\centering
\caption{Background framework configuration parameters.}
\label{tab:background-config}
\small
\begin{tabular}{lll}
\hline
\textbf{Parameter} & \textbf{Default} & \textbf{Description} \\
\hline
\texttt{BACKGROUND\_ENABLED} & true & Master enable/disable switch \\
\texttt{BACKGROUND\_HEALTH\_ENABLED} & true & Enable health monitoring \\
\texttt{BACKGROUND\_HEALTH\_INTERVAL\_SECONDS} & 30 & Health check interval \\
\texttt{BACKGROUND\_BACKUPS\_ENABLED} & false & Enable scheduled backups \\
\texttt{BACKGROUND\_BACKUPS\_CRON} & \texttt{30 2 * * *} & Backup cron schedule (UTC) \\
\texttt{BACKGROUND\_CLEANUP\_ENABLED} & false & Enable cleanup tasks \\
\texttt{BACKGROUND\_CLEANUP\_CRON} & \texttt{15 3 * * 0} & Cleanup cron (UTC) \\
\texttt{BACKGROUND\_REDIS\_CLEANUP\_ENABLED} & true & Enable Redis index cleanup \\
\texttt{BACKGROUND\_REDIS\_CLEANUP\_INTERVAL} & 3600 & Cleanup interval (seconds) \\
\texttt{BACKGROUND\_REDIS\_CLEANUP\_BATCH\_SIZE} & 500 & Members per cleanup batch \\
\hline
\end{tabular}
\end{table}

